\documentclass[11pt]{article}
\usepackage[margin=1in]{geometry}
\usepackage[T1]{fontenc}
\usepackage{lmodern,microtype,amsmath,amssymb,amsthm}
\usepackage[colorlinks=true,linkcolor=blue,citecolor=blue,urlcolor=blue]{hyperref}
\hypersetup{pdftitle={RH-Conditional Approximation Estimates and the Remaining Simplicity Problem},pdfauthor={Drafted for Edward Baker}}
\usepackage{enumitem}
\setlist{nosep,leftmargin=*}
\newtheorem{theorem}{Theorem}
\newtheorem{proposition}[theorem]{Proposition}
\DeclareMathOperator{\Log}{Log}
\newcommand{\RH}{\mathrm{RH}}
\newcommand{\E}{\mathcal E}
\newcommand{\dd}{\,\mathrm d}
\newcommand{\Nref}{\href{../notes/16_GEVREY_GLOBAL_APPROXIMATION_AND_RATE_20261020.md}{Note 16}}
\newcommand{\Href}{\href{../notes/15_COMPLETED_RATIO_ENERGY_AND_MULTIPLICITY_20261020.md}{Note 15}}
\newcommand{\CNref}{\href{../notes/17_NORMALIZED_HEAT_GLOBAL_LIMIT_AND_ARITHMETIC_COSTS_20261020.md}{Note 17}}
\newcommand{\CTref}{\href{../notes/18_ARITHMETIC_HEAT_TRANSFER_AND_EXACT_COMPLETION_20261021.md}{Note 18}}
\title{RH-Conditional Approximation Estimates\\and the Remaining Simplicity Problem}
\author{Drafted for Edward Baker}
\date{21 October 2026\\\small Research consolidation through Notes 17--18}
\begin{document}
\maketitle
\begin{abstract}
We consolidate analytic estimates obtained while investigating whether the
Riemann hypothesis implies simplicity of the nontrivial zeros of the zeta
function. The principal approximation argument establishes, within the
research notes, a complete first-order limit for fixed Gevrey-tapered
M\"obius polynomials and an RH-conditional infimum constant equal to the
sum of squared multiplicities. A completed-ratio energy has an explicit
heat-scale transition governed by the same multiplicity excess. Exact
finite Dirichlet filters give a sharp local resolution cost. A normalized
heat target admits a complete first-order limit at bounded scaled time.
We compute the sharp cost of transferring its profiles back to the
original arithmetic problem over fixed finite filtered-Gevrey families,
and show that specified exact heat completions still fail to retain an
ordinary Dirichlet series at fixed positive time. These estimates do
not exclude a multiple zero: the missing ingredient is an
arithmetic upper estimate that contradicts its additional cost. This
brief presents precise statements and proof mechanisms, with the full
analytic payments referenced to the research notes.
\end{abstract}

\section{Scope and the unresolved quantity}
Unless stated otherwise, RH is an assumption, not a conclusion. Write
\(\rho=\frac12+i\gamma\) for a distinct nontrivial zero, \(m_\rho\) for
its multiplicity, and sum over both ordinate signs. The two masses are
\begin{equation}\label{eq:KZ}
 K=\sum_{\rho\ {\rm distinct}}\frac{m_\rho^2}{|\rho|^2},
 \qquad
 Z=\sum_{\rho\ {\rm distinct}}\frac{m_\rho}{|\rho|^2}
   =2+\gamma_E-\log(4\pi).
\end{equation}
They are finite; the usual unit-window zero count also bounds the sum
of squared multiplicities with this weight. Their difference is
\begin{equation}\label{eq:gap}
 K-Z=\sum_{\rho\ {\rm distinct}}
            \frac{m_\rho(m_\rho-1)}{|\rho|^2}.
\end{equation}
A double conjugate pair contributes \(4/(\gamma^2+1/4)\).
The approximation and heat estimates below retain this excess.
Identifying it is a bookkeeping conclusion, not a new simplicity
criterion or a proof that the excess vanishes.

The results inherited from boundary winding justify the local
multiplicity calculation and finite-height continuation under RH.
They do not furnish a strict surplus with the opposing sign.
The present judgment is therefore that the investigation has produced
analytic estimates, but does not have a nearly complete simplicity
proof.

\paragraph{Status.}
The new analytic results are proved in the linked research notes and
have undergone separate same-model internal derivations and adversarial
checks. They have not received specialist external review or analytic
formalization. In particular, the complete scaled-tail theorem below is
a substantial internal premise, not a previously published theorem
being cited under a new name. No novelty or priority claim is made.

\section{Fixed Gevrey approximation and complete tails}
For a Dirichlet polynomial \(D\), define
\[
 E(D)=\frac1{2\pi}\int_{\mathbb R}
  \frac{|1-\zeta(\frac12+it)D(\frac12+it)|^2}{1/4+t^2}\dd t,
 \qquad
 d_N^2=\inf_{D(s)=\sum_{n\le N}a_n n^{-s}}E(D).
\]
We also use \(\|g\|_{\mathcal H}^2=(2\pi)^{-1}
\int_{\mathbb R}|g(\frac12+iy)|^2(1/4+y^2)^{-1}\dd y\),
so \(E(D)=\|1-\zeta D\|_{\mathcal H}^2\).
Let \(W\) be real, smooth, supported in \([-1,1]\), with \(W(0)=1\)
and
\[
 \|W^{(j)}\|_1\le A^{j+1}(j!)^\sigma,\qquad 1<\sigma<2.
\]
Set \(w=W|_{[0,1]}\), \(L=\log N\), and
\[
 V_{N,w}(s)=\sum_{n\le N}\mu(n)n^{-s}w(\log n/L),\qquad
 p_{m,w}(u)=-\frac{(u^{m-1}w(u))^{(m)}}{(m-1)!},\qquad
 C_m(w)=\|p_{m,w}\|_{L^2(0,1)}^2.
\]
\begin{theorem}[Complete fixed-taper estimate]\label{thm:ST}
Assume RH and fix \(W\) as above. If \(E_{N,w}^{>T}\) denotes the
part of \(E(V_{N,w})\) at \(|t|>T\), then for every
\(0<\eta<2-\sigma\),
\begin{align}
 \sup_{N\ge2}(\log N)E_{N,w}^{>T}
      &\ll_{W,\eta}T^{\sigma-2+\eta},\qquad T\ge2,
       \label{eq:tail}\\
 \lim_{N\to\infty}(\log N)E(V_{N,w})
      &=\sum_{\rho\ {\rm distinct}}
                \frac{C_{m_\rho}(w)}{|\rho|^2}<\infty.
       \label{eq:fixed}
\end{align}
No simplicity, zero-gap lower bound, or inverse-derivative moment is
assumed.
\end{theorem}

\paragraph{Proof mechanism and dependencies.}
The full proof is \Nref, Sections 2--6, 8 and 11.
For \(s=\frac12+it\) and \(H>0\), its exact starting identity is
\[
 V_{N,w}(s)=\int_{\mathbb R}
 \frac{K_H(v)}{\zeta(s+(H+iv)/L)}\dd v,\quad
 K_H(v)=\frac1{2\pi}\int_{\mathbb R}W(u)e^{(H+iv)u}\dd u .
\]
Fourier inversion on an absolutely convergent M\"obius line, followed
by a shift within the RH zero-free half-plane, preserves every
coefficient. Integration by parts against the whole complex
exponential gives
\[
 |K_H(v)|\le C_W e^H e^{-c_W|v|^{1/\sigma}},
 \qquad \int_{\mathbb R}K_H(v)\dd v=1 ,
\]
with constants independent of \(H>0\). The factor \(e^H\) is retained.

The primary height inputs are the RH argument bound of
Carneiro--Chandee--Milinovich~\cite{CCM} and the subpower zeta bound of
Chandee--Soundararajan~\cite{CS}. In the low-height range, the proof
keeps the whole nearby zero cluster, including repetitions. Its
kernel moment has logarithm
\((\sigma-1)M\log(M+2)+O_W(M+1)\) for a cluster of size \(M\).
This is paid using the half-coefficient local counting allowance;
no individual reciprocal derivative is estimated. A separate moving
line handles the high-height range and its complete remote convolution.

Subtracting the constant term coherently gives decay in the actual
distance to the nearest zero. Squaring and integrating that decay over
nearest-zero cells costs \(O(1/L)\) independently of the gaps.
This proves \eqref{eq:tail}. On a fixed isolated zero neighborhood,
the complete error has profile
\(\int_0^1p_{m,w}(u)e^{-iLyu}\dd u+O_{W,\rho}(L^{-1})\),
with constants allowed to depend on that neighborhood. Its unscaled
norm is \(C_m(w)/(L|\rho|^2)+O(\log L/L^2)\).
The separately established tail bound then permits height exhaustion.
Thus local asymptotics alone are not being used as a global upper bound.

\paragraph{A controlled moving family.}
For the almost-linear taper
\(w_\eta(u)=(1-u)\chi((1-u)/\eta)\), with a fixed Gevrey cutoff
\(\chi=0\) on \((-\infty,\frac12]\) and \(\chi=1\) on \([1,\infty)\),
a separate proof gives
\begin{equation}\label{eq:moving}
 E(V_{N,w_\eta})\le \frac{C\exp(C\eta^{-a})}{L},
 \quad L^{-\theta_0}\le\eta<\tfrac14,\quad
 a>\frac1{2-\sigma},\quad 0<\theta_0<1.
\end{equation}
Here the parameters \(\sigma,\theta_0,a\) are fixed
(\Nref, Section 12). Consequently
\(\eta_N=(\log L)^{-\kappa}\), \(0<\kappa<2-\sigma\),
gives \(E(V_{N,w_{\eta_N}})\le L^{-1+o(1)}\).
The exponential allowance is compatible with the multiple-zero lower
costs. It supplies neither a moving first-order asymptotic nor a
multiplicity-free upper constant.

\section{The infimum constant and finite Dirichlet filters}
Integration by parts gives
\(\int_0^1u^jp_{m,w}(u)\dd u=\delta_{j0}\) for \(0\le j<m\).
Projection onto polynomials of degree below \(m\) therefore yields
\begin{equation}\label{eq:localfloor}
 C_m(w)=m^2+\|p_{m,w}-\mathcal K_m(\,\cdot\,,0)\|_2^2\ge m^2,
\end{equation}
where \(\mathcal K_m\) is their endpoint reproducing kernel on \([0,1]\).
Flat Gevrey approximations to \((1-u)^m\) approach this lower cost.
One common taper need not optimize different multiplicities at once.

\begin{theorem}[Infimum consequence of Theorem~\ref{thm:ST}]
Assuming RH and the complete fixed-taper estimate,
\begin{equation}\label{eq:inf}
             \lim_{N\to\infty}(\log N)d_N^2=K.
\end{equation}
Real coefficients suffice.
\end{theorem}
\begin{proof}[Proof outline]
Burnol's lower bound~\cite{Burnol} gives the lower limit \(K\).
For the upper bound, fix a baseline taper and a finite set \(F\) of
zeros carrying all but a small part of its profile mass. Choose a
fixed near-optimal taper separately at each target.
Bounded finite Dirichlet multipliers \(P_\rho\) select these profiles:
\[
 D_N=V_{n,w_0}+\sum_{\rho\in F}P_\rho
                  (V_{n,w_\rho}-V_{n,w_0}),\qquad n=\lfloor N/M\rfloor,
\]
where all multipliers have length at most \(M\).
The multiplier coefficients of the error sum identically to one,
so every cross term and moment constraint is retained.
For fixed multipliers, the local profile expansion together with
Theorem~\ref{thm:ST}'s scaled-tail bound gives the complete mixed limit. Its profile tail is bounded by a fixed multiple of the sum
of the finitely many original profile tails.

One may interpolate exactly on any larger finite zero set \(G\).
For \(0<\theta<\pi\), \(k=\lfloor\pi/\theta\rfloor\), the polynomial
\[
 B_\theta(z)=
 \frac{\sin\theta-\sin((k+1)\theta)z^k+\sin(k\theta)z^{k+1}}
      {\sin\theta-\sin((k+1)\theta)+\sin(k\theta)}
\]
has nonnegative coefficients summing to one, equals one at \(1\),
and vanishes at \(e^{\pm i\theta}\). At \(\theta=\pi\), use
\((1+z)/2\). Substitute \(z=p^{\rho-s}\), choosing \(p=2\) or \(3\)
with nonzero relative phase, and multiply over the other nodes of
\(G\). Irrationality of \(\log2/\log3\) guarantees such a choice.
The products remain bounded by one on the critical line.
Their complete exterior profile allowance is independent of their
degrees and tends to zero as \(G\) exhausts the zeros.

Keep \(F\), tapers and finite \(G\) fixed before taking \(N\to\infty\).
Then exhaust \(G\), optimize the finite profile costs, and let the
initial baseline-tail tolerance tend to zero. Conjugate
pairing makes the final coefficients real. This proves the upper
limit \(K\); see \Nref, Sections 13--14.
\end{proof}

This conclusion concerns the unrestricted polynomial infimum, not
the original linear M\"obius polynomial. It removes a common-taper
restriction but leaves the squared multiplicities in \eqref{eq:KZ}.
It provides no effective threshold in \(N\).

\begin{proposition}[Sharp two-point isolation cost]\label{prop:filter}
Let \(P\) have Dirichlet length at most \(M\), be bounded by one on
\(\Re s=\frac12\), and satisfy \(P(\rho)=1\) and
\(|P(\rho+i\Delta)|\le\varepsilon\), \(0\le\varepsilon<1\).
Then
\[
 |\Delta|\log M\ge2\arccos\varepsilon.
\]
For exact cancellation at a small nonzero gap, the least log-length
lies between \(\pi/|\Delta|\) and \(\pi/|\Delta|+\log2\).
\end{proposition}
\begin{proof}
Centering the frequency interval gives an entire function of type
at most \((\log M)/2\). Its real part starts at one, is bounded by one,
and ends at most \(\varepsilon\).
The sharpened Bernstein inequality~\cite{Redheffer}, integrated after
arccos, gives the lower bound. The polynomial \(B_\theta\) with
\(p=2\) and \(\theta=|\Delta|\log2\) gives the upper bound.
\end{proof}
This prices root isolation only. Assigning the same taper to two
roots need not isolate them, and the proposition is not a lower
bound for every approximation scheme.

\section{Completed-ratio energy and heat separation}
Use the original theta source
\[
 H(t,x)=\int_0^\infty e^{tu^2}\Phi(u)\cos(xu)\dd u,\qquad
 \Phi(u)=\sum_{n\ge1}(2\pi^2n^4e^{9u}-3\pi n^2e^{5u})
                                      e^{-\pi n^2e^{4u}},
\]
with \(\xi(s)=\tfrac12s(s-1)\pi^{-s/2}\Gamma(s/2)\zeta(s)\).
The completion and heat normalization are
\[
 \mathcal F_t(z)=8H(t,-2iz),\quad \mathcal F_0(z)=\xi(\tfrac12+z),\quad
 \partial_t\mathcal F_t=\tfrac14\partial_z^2\mathcal F_t.
\]
For \(\alpha>0\), define
\[
 J_t(\epsilon)=\frac1{2\pi}\int_{\mathbb R}
 \frac{|1-\mathcal F_t(iy)/\mathcal F_t(\epsilon+iy)|^2}
      {\alpha^2+y^2}\dd y,\qquad
 M_\alpha=\sum_\gamma\frac{m_\gamma}{\alpha^2+\gamma^2},\quad
 K_\alpha=\sum_\gamma\frac{m_\gamma^2}{\alpha^2+\gamma^2}.
\]
Under RH, for \(t\ge0\), the paired order-one product makes the ratio contractive.
Its positive harmonic majorant has a known complete mass.
Local profiles followed by exhaustion of that mass give
\begin{equation}\label{eq:ratio}
 \lim_{\epsilon\downarrow0}\frac{J_0(\epsilon)}{\epsilon}
    =\sum_\gamma\frac{c_{m_\gamma}}{\alpha^2+\gamma^2},
 \qquad c_m=m\left(1-\frac{\binom{2m}{m}}{4^m}\right).
\end{equation}
In particular \(c_1=\frac12\), \(c_2=\frac54\); no mass at infinity
is discarded. The proof is \Href, Sections 1--3.

For \(t=v\epsilon^2\), the local heat expansion at an order-\(m\)
zero gives Hermite centers \(r\lambda_j\), where
\(r=\sqrt{v/2}\) and \(\mathrm{He}_m(\lambda_j)=0\).
Define
\[
 P_{m,r}(y)=\prod_{j=1}^m(y-r\lambda_j),\qquad
 \mathcal C_m(r)=\frac1{2\pi}\int_{\mathbb R}
 \left|1-\frac{P_{m,r}(y)}{P_{m,r}(y-i)}\right|^2\dd y.
\]
The same complete-mass argument proves, at each fixed \(v\ge0\),
\[
 \lim_{\epsilon\downarrow0}\frac{J_{v\epsilon^2}(\epsilon)}{\epsilon}
   =\mathcal A_\alpha(v)
   :=\sum_\gamma
      \frac{\mathcal C_{m_\gamma}(\sqrt{v/2})}{\alpha^2+\gamma^2}.
\]
The analytic transition is quantified by
\begin{align}
 \mathcal C_2(r)&=1+\frac1{4(1+r^2)},\label{eq:double}\\
 0\le \mathcal A_\alpha(v)-\frac{M_\alpha}{2}
     &\le\frac{K_\alpha-M_\alpha}{4v},\qquad v>0,\label{eq:heatbound}\\
 \lim_{v\to\infty}v\left(\mathcal A_\alpha(v)-\frac{M_\alpha}{2}\right)
     &=\frac{K_\alpha-M_\alpha}{4}.\label{eq:heatlimit}
\end{align}
The estimate follows from pairwise overlap integrals of the local
contractive factors and
\(\sum_{i<j}(\lambda_i-\lambda_j)^{-2}=m(m-1)/8\).
The weighted multiplicity-square sum dominates the passage to all
zeros. At \(\alpha=\frac12\), the last coefficient is \((K-Z)/4\).

These are ordered limits: first \(\epsilon\to0\) at fixed \(v\), then
\(v\to\infty\). They do not assert convergence of time derivatives
or a joint limit with unbounded \(v(\epsilon)\).
There is also a global obstruction to using the original completion
\(A(s)=\tfrac12s(s-1)\pi^{-s/2}\Gamma(s/2)\).
For every fixed \(0<t\le\frac12\) and nonzero finite Dirichlet polynomial
\(D\), the original weighted error against
\(Z_t(s)=\mathcal F_t(s-\frac12)/A(s)\) is infinite, unconditionally.
Polymath's two-carrier asymptotic~\cite{Polymath}, with finite-frequency
integration by parts, gives a block lower bound
\(T^{-1}\exp[(t/8)\log^2(T/(2\pi))]\to\infty\).
For \(b\ge0\), multiplying the weight at large heights by
\(e^{-b\log^2(|y|/(2\pi))}\) restores finiteness exactly when
\(b\ge t/8\) (\Nref, Section 16). This changes the norm; it leaves
the compact heat calculation and completed-ratio energy intact.

\section{A complete normalized heat limit}\label{sec:normalized}
The divergence just described motivates removing the leading completion
carrier while keeping the original polynomial and boundary weight. Put
\begin{equation}\label{eq:normalized-target}
 a(s)=\frac1{2s}+\frac1{s-1}+\frac12\Log\frac{s}{2\pi},\qquad
 T_\tau(s)=e^{-\tau a(s)^2/4}
                \frac{\mathcal F_\tau(s-\frac12)}{A(s)}.
\end{equation}
The principal logarithm is used on the critical line. For fixed
\(v\ge0\), set
\[
 E_{N,w}^{[v]}=\|1-T_{v/L^2}V_{N,w}\|_{\mathcal H}^2,\qquad
 f_m(u)=\frac{u^{m-1}w(u)}{(m-1)!},
\]
and define the heat profile
\begin{equation}\label{eq:heat-profile}
 h_{m,v,w}=-\sum_{k=0}^{\lfloor m/2\rfloor}
       \frac{m!}{k!(m-2k)!}\left(\frac v4\right)^k f_m^{(m-2k)}.
\end{equation}
Here \(h\) denotes an approximation profile, not a physical derivative
of \(H\).

\begin{theorem}[Normalized complete limit]\label{thm:CN}
Assume RH and fix the normalized Gevrey taper of
Theorem~\ref{thm:ST}. For every fixed \(V<\infty\) and
\(0<\eta<2-\sigma\),
\begin{align}
 \sup_{N\ge2}\sup_{0\le v\le V}
       L(E_{N,w}^{[v]})^{>T}
       &\ll_{W,V,\eta}T^{\sigma-2+\eta},\qquad T\ge2,
       \label{eq:CN-tail}\\
 \lim_{N\to\infty}L E_{N,w}^{[v]}
       &=\sum_{\rho\ {\rm distinct}}
          \frac{\|h_{m_\rho,v,w}\|_2^2}{|\rho|^2}<\infty.
       \label{eq:CN-limit}
\end{align}
The limit is locally uniform in \(v\ge0\). The taper is fixed;
no joint limit with unbounded \(v=v_N\) is asserted.
\end{theorem}

\paragraph{Proof mechanism.}
\CNref, Sections 1--5, starts with the exact complete-source identity
\[
 \mathcal F_\tau(s-\tfrac12)=\frac1{\sqrt\pi}
       \int_{\mathbb R}e^{-x^2}\xi(s+\sqrt\tau x)\dd x.
\]
Gaussian domination justifies both integrations before any approximation.
Carrier cancellation and reflection give a subpower bound for
\(T_\tau\) uniformly over bounded \(\tau\ge0\). Combining this
Gaussian integral with the M\"obius convolution retains both remote
tails and the entire nearby zero cluster. The same nearest-zero-cell
argument as in Theorem~\ref{thm:ST} pays the scaled height tail.
The compact heat polynomial at an order-\(m\) zero then gives
\eqref{eq:heat-profile}; only after \eqref{eq:CN-tail} is established
is the compact calculation exhausted over all heights.

At \(m=1\) the profile is unchanged. At \(m=2\) it is
\(h_{2,v,w}=-(uw)''-(v/2)uw\). In particular the squared local norm
of \((T_{v/L^2}-\zeta)V_{N,w}\) at an exact double zero is
\[
 \frac{v^2\|uw\|_2^2}{4L|\rho|^2}+O_{v,w,\rho}(L^{-2});
\]
the norm here is over a fixed isolated neighborhood of that zero.
It occurs at the same first-order scale as the proposed heat gain.

\paragraph{Analytic and arithmetic limitations.}
For each fixed \(t>0\), \CNref, Sections 6--7, proves unconditionally,
with \(\ell=\log(\sigma/(2\pi))\) and real \(\sigma\to+\infty\),
\begin{equation}\label{eq:carrier-defect}
 T_t(\sigma)=1+\frac{t^2\ell^2/64+t/8}{\sigma}
             +O_t\!\left(\frac{(1+|\ell|)^4}{\sigma^2}\right).
\end{equation}
An ordinary Dirichlet series converging in a right half-plane has
remainder \(O(2^{-\sigma})\) after its constant term, so cannot
represent this target, even if its initially assumed convergence is
conditional. Separately, \(s=1\) becomes an essential singularity:
\[
 T_t(1+z)=z^{-1}e^{-t/(4z^2)-tb_0/(2z)}U_t(z),\qquad
 b_0=\frac{1-\log(2\pi)}2,
\]
where \(U_t\) is holomorphic and \(U_t(0)>0\). A finite-order zero
or pole subtraction cannot repair it. The essential singularity alone
would not exclude a Dirichlet series farther right; that conclusion
uses \eqref{eq:carrier-defect}. Thus the finite critical-line norm does
not restore the original Nyman--Beurling/Hardy hypotheses or authorize
the original arithmetic lower bound for this changed target.

\section{The sharp arithmetic repair of a heat profile}\label{sec:repair}
Heating the finite coefficients instead of the numerator is exact:
\begin{equation}\label{eq:coefficient-heat}
 e^{\frac{v}{4L^2}\partial_s^2}V_{N,w}=V_{N,w_v},\qquad
 w_v(u)=e^{vu^2/4}w(u).
\end{equation}
This stays within the fixed Gevrey class, but its original-arithmetic
profile is \(q_{m,v,w}=-(e^{vu^2/4}f_m)^{(m)}\), not
\(h_{m,v,w}\). In fact
\[
 \left.\partial_v(q_{m,v,w}-h_{m,v,w})\right|_{v=0}
 =-\tfrac14\{u^2f_m^{(m)}+2mu f_m^{(m-1)}\}\ne0.
\]
The braces equal \(2mu+O(u^2)\) near zero. Even a simple zero has
this coefficient-transfer defect. The two complete tail theorems give
the actual all-zero transfer limit by polarization; no infinite
heat-evolved Euler series is used.

Let \(\Pi_m\) project onto the polynomials of degree below \(m\),
and abbreviate \(K_m=\mathcal K_m(\,\cdot\,,0)\). Every original
arithmetic profile satisfies \(\Pi_m p=K_m\), by its first \(m\)
moments. The least squared correction to a specified heat profile is
therefore
\begin{equation}\label{eq:repair}
 \delta_{m,v}(w)^2=\|K_m-\Pi_m h_{m,v,w}\|_2^2,\qquad
 p_*=h_{m,v,w}+K_m-\Pi_m h_{m,v,w}.
\end{equation}
This distance is sharp in the closure of real normalized flat Gevrey
tapers, without a positivity or monotonicity constraint. To see this,
solve \(f^{(m)}=-p_*\) with the normalized initial jets. The moment
conditions force \(f^{(j)}(1)=0\) for \(j<m\); flattening \(f\)
near one changes its \(m\)th derivative by \(O(\eta^{1/2})\)
in \(L^2\). Orthogonality makes the excess squared correction
\(O(\eta)\). See \CTref, Sections 1--3.

The correction is explicit. If
\(b_j=\int_0^1u^j(p_{m,w}-h_{m,v,w})\dd u\) and
\(G_m=(1/(i+j+1))_{0\le i,j<m}\), then
\[
 \delta_{m,v}(w)^2=b^*G_m^{-1}b
       \ge (m-\|h_{m,v,w}\|_2)_+^2.
\]
For a double zero, with \(I_j=\int_0^1u^j(uw(u))\dd u\),
\begin{equation}\label{eq:double-repair}
 \delta_{2,v}(w)^2=\frac{v^2}{4}
                 (4I_0^2-12I_0I_1+12I_1^2).
\end{equation}
A double heat profile with squared norm below two consequently requires
a correction strictly greater than \((2-\sqrt2)^2\) at that zero,
before division by \(|\rho|^2\).

\begin{theorem}[Complete optimal transfer in the declared class]
\label{thm:transfer}
Assume RH and the complete estimates of Theorems~\ref{thm:ST}
and~\ref{thm:CN}. Fix \(v,w\). Consider families
\[
 D_N=\sum_{j=0}^J A_j(s)V_{\lfloor N/M\rfloor,w_j}(s),\qquad
 \sum_{j=0}^J A_j\equiv1,
\]
with fixed finite Dirichlet multipliers supported through \(M\) and
fixed normalized Gevrey tapers. Put
\(q_\rho=\sum_jA_j(\rho)p_{m_\rho,w_j}\). Then
\begin{align}
 \mathcal T(D;v,w)&:=\lim_{N\to\infty}L
     \|\zeta D_N-T_{v/L^2}V_{N,w}\|_{\mathcal H}^2
       =\sum_\rho\frac{\|q_\rho-h_{m_\rho,v,w}\|_2^2}{|\rho|^2},
       \label{eq:transfer-limit}\\
 \inf_D\mathcal T(D;v,w)&=
       \Delta_v(w):=\sum_\rho\frac{\delta_{m_\rho,v}(w)^2}{|\rho|^2}
       <\infty.\label{eq:transfer-inf}
\end{align}
The infimum is over these fixed families after taking their \(N\)
limits, not over arbitrary coefficient sequences depending on \(N\).
\end{theorem}
\begin{proof}[Proof outline]
The compact mixed profiles and the two complete tail estimates prove
\eqref{eq:transfer-limit}. Projection gives the lower bound in
\eqref{eq:transfer-inf}. Use the original taper as a baseline with
finite total transfer cost, repair a finite set of profiles to within
any tolerance, and assign them using the bounded two-prime filters.
The complete majorant includes the baseline and every selected taper
and is independent of filter degree. Dominated convergence pays the
exterior profile mass. First take \(N\to\infty\) for each fixed
family, then increase filter degree and optimize tapers. This gives
the matching upper bound (\CTref, Section 5).
\end{proof}

\paragraph{An exact example distinguishing correction from final cost.}
At \(v=2\pi^2\), take the order-two closure profile
\[
 f(u)=\frac{(1-u)\sin(\pi u)}\pi,\qquad
 w(u)=f(u)/u,\qquad h=-f''-\pi^2f=2\cos(\pi u).
\]
Its heat cost is \(\|h\|_2^2=2\), whereas
\begin{align*}
 \delta^2&=4-\frac{48}{\pi^2}+\frac{192}{\pi^4},\\
 p_*&=2\cos(\pi u)+4-\frac{24}{\pi^2}
                   +\left(-6+\frac{48}{\pi^2}\right)u,\\
 \|p_*\|_2^2&=6-\frac{192}{\pi^4}>4.
\end{align*}
Fixed flat cutoffs realize these local limits, with \(N\to\infty\)
before the cutoff width decreases. The closest repair to \(h\) need
not minimize the repaired profile's own norm. This example does not
raise the original infimum above \(K\); it quantifies the price of
retaining this particular heat profile. No uniform full-height limit
for a moving base cutoff is claimed.

\section{Exact completion still changes the arithmetic}\label{sec:completion}
The algebraic error in \eqref{eq:carrier-defect} might suggest heating
the completion exactly. \CTref, Sections 6--10, treats that alternative
without assuming RH. Write \(X_t(s)=\mathcal F_t(s-\frac12)\).
For fixed \(t>0\) and \(c>1\), define the entire first-mode factor
\begin{equation}\label{eq:exact-completion}
 B_t^c(s)=\frac1{\sqrt{\pi t}}\int_c^\infty
                     e^{-(x-s)^2/t}A(x)\dd x.
\end{equation}
It is positive on the real axis and solves
\(\partial_tB_t^c=\frac14\partial_s^2B_t^c\).
Splitting the Gaussian formula for \(X_t\) at \(c\), and expanding
the absolutely convergent Euler series only on \(x\ge c\), gives
\begin{equation}\label{eq:mode-decomposition}
 X_t(s)=Q_t^c(s)+\sum_{n\ge1}n^{-s}e^{t(\log n)^2/4}
                     B_t^c(s-\tfrac t2\log n).
\end{equation}
The series converges absolutely and locally uniformly for every complex
\(s\); its displayed coefficients depend on \(s\). The entire
left contribution satisfies
\(Q_t^c(\sigma)=O_{t,c}(e^{-(\sigma-c)^2/t})\) for real
\(\sigma\ge c\). This retains the full negative-real-axis integral.

\begin{proposition}[Fixed-time second-mode obstruction]
\label{prop:completion}
For fixed \(t>0,c>1\), as real \(\sigma\to+\infty\),
\begin{equation}\label{eq:second-mode}
 \frac{X_t(\sigma)}{B_t^c(\sigma)}-1
 =2^{-\sigma}e^{t(\log2)^2/4}
    \left(\frac{\sigma}{2\pi}\right)^{-t\log2/4}
    \left[1+O_{t,c}\!\left(\frac{\log\sigma}{\sigma}\right)\right].
\end{equation}
This quotient has no ordinary Dirichlet-series representation converging
in a right half-plane, even conditionally.
\end{proposition}
\begin{proof}[Proof mechanism]
A fixed-shift Laplace estimate for \(B_t^c\) gives the \(n=2\)
term. Monotonicity on the real axis bounds the complete \(n\ge3\)
tail by \(O_{t,c}(3^{-\sigma}\sigma^{-t\log3/4})\); the left
term is Gaussian-small. A hypothetical Dirichlet series has expansion
\(d_1+d_2 2^{-\sigma}+O(3^{-\sigma})\) sufficiently far right.
Equation~\eqref{eq:second-mode} forces \(d_1=1,d_2=0\), but its
nonzero second-mode remainder is not \(O(3^{-\sigma})\).
\end{proof}

The same conclusion holds for the explicit contour prescriptions
\[
 B_{t,\eta}(s)=\frac1{\sqrt{\pi t}}
     \int_{\mathbb R+i\eta}e^{-(z-s)^2/t}A(z)\dd z,
       \qquad \eta\in\mathbb R\setminus\{0\},
\]
and their symmetric upper/lower average. Gamma reflection pays the
negative horizontal ray. Deforming only the pole-free right part gives
\(B_{t,\eta}(\sigma)-B_t^c(\sigma)
=O_{t,c,\eta}(e^{-(\sigma-c)^2/t})\), preserving
\eqref{eq:second-mode}. More generally the comparison applies to
analytic normalizers differing from \(B_t^c\) by a Gaussian-small
term on this ray. It is not an exclusion of arbitrary completions.

Wherever a heat-solution normalizer \(B\) is nonzero, its quotient
obeys
\[
 \partial_t(X_t/B)=\tfrac14\partial_s^2(X_t/B)
               +\tfrac12(B'/B)\partial_s(X_t/B).
\]
Thus exact division leaves a completion drift. No global zero-free
property, critical-line norm theorem for these exact quotients, or
uniform assertion for \(t=t_N\to0\) is established here. In particular,
Theorem~\ref{thm:CN} is not being reassigned to these new normalizers.

\section{Relation to BW09}
The originating BW09 route seeks to rule out a hypothetical positive
all-real threshold by forcing a later collision. Its current manuscript
assumes an attained finite-height ordinary event at that threshold;
coverage of every hypothetical positive threshold remains a separate
unproved step. BW10 assumes RH and starts from a multiple zero at the
specified time zero, allowing arbitrary finite multiplicity. It removes
that unknown-time and attainment issue from its conditional problem.
This is a structural distinction, not a proved ordering of difficulty.

The unchanged theta source, local Hermite splitting, root-wall
cancellation and complete boundary-phase identities are shared.
The ratio-energy analysis also applies relative to a known all-real
base time \(T\), with source \(e^{Tu^2}\Phi(u)\), provided the
all-real forward continuation and order-one product hypotheses are
retained. The finite-selector estimates are unconditional filter lemmas.
These tools sharpen the accounting of an existing event; they do not
supply additional winding.

The strongest approximation estimates have a narrower scope.
Theorem~\ref{thm:ST} uses RH for the original zeta function and its
unchanged M\"obius convolution. Importing that assumption into BW09's
positive-threshold contradiction would assume the desired conclusion.
All-realness of a positive-time heat deformation supplies neither this
original Dirichlet arithmetic nor a substitute theorem.
The normalized heat and sharp transfer theorems inherit this RH premise.
In contrast, the fixed-positive-time completion obstructions in
Sections~\ref{sec:normalized} and~\ref{sec:completion} are unconditional
and may inform either route, within their stated normalization classes.

BW09 still needs a complete theta estimate forcing an ordinary born
pair's cap-phase decrease strictly beyond \(\pi\), with wall
continuation, a regular terminal cap, endpoint deficits and source
remainders all controlled. Its fixed-index large-time cosine model
produces the ordinary \(\pi\), not a surplus. BW10 has not supplied
that strict estimate or the threshold-event coverage step. See the
\href{../../09_expanding_contour_later_collision/manuscript/main.tex}{BW09 manuscript},
Sections on scope, root walls, surplus and information limits, and its
\href{../../09_expanding_contour_later_collision/research-progress/03_ANALYTICAL_PHASE_SURPLUS_PROGRESS_20261020.md}{phase-surplus progress register}.

\section{What is missing, and what to brainstorm next}
The common bottleneck is now explicit. The complete approximation
upper bound has coefficient \(K\), and the heat transition has
excess coefficient \((K-Z)/4\). Both are compatible with a multiple
endpoint. A simplicity proof needs an additional estimate using the
unchanged arithmetic of zeta that rules out this excess, or a different
mechanism that excludes an original multiple zero.
Sections~\ref{sec:normalized}--\ref{sec:completion} make the transfer
issue quantitative: normalized numerator heat has a complete norm,
but returning its profiles to the original arithmetic class costs
\(\Delta_v(w)\), and the tested exact completions retain a
non-Dirichlet second mode. These results close specific unpaid
transfers; they do not rule out every use of heat flow or every
possible completion.

The tested absolute prime and M\"obius-correlation payments do not
supply that estimate. In the exact ratio-energy representation, signed
off-diagonal terms cancel a diagonal of size \(1/\epsilon\) to bounded
size, with the residual energy determined at absolute precision
\(O(\epsilon)\); taking absolute values gives a divergent
double payment (\Href, Section 6). The prime-error linearization
also has a first-order remainder at multiple zeros (Section 7).
Neither remainder can be omitted at the scale where simplicity is
being tested. Generic positive-source collision controls demonstrate
the insufficiency of source positivity alone, while failing the
ordinary Dirichlet-series premise; they are not counterexamples
satisfying all of zeta's arithmetic properties.

The next discussion should focus on two questions:
\begin{enumerate}
\item Is there a specific signed arithmetic identity whose \emph{complete}
remainder can be bounded more sharply than these multiplicity-compatible
estimates? A proposed bound must retain all zero and prime tails.
\item Is there a theta-specific sign or rigidity principle, absent from
the screened collision controls, that excludes the multiple endpoint
without merely rewriting the desired conclusion?
\end{enumerate}
Before further expansion, the complete scaled-tail theorem deserves
external scrutiny because several later conclusions depend on it.
The present estimates warrant consolidation and a change of viewpoint,
not a claim of proximity to a proof.

\paragraph{Supporting checks.}
The \href{../numerics/EXACT_ALGEBRA_REPLAY_20261020.json}{exact-algebra replay}
and \href{../numerics/lean/README.md}{bounded Lean check}
support selected identities and an inner-product tradeoff.
The \href{../../09_expanding_contour_later_collision/manuscript/main.tex}{BW09 manuscript}
cites supporting numerics. These checks do not verify the analytic tail
theorem or simplicity and are not proof premises here.

\paragraph{Acknowledgement and draft record.}
Drafted for Edward Baker with substantial OpenAI \texttt{gpt-6-astra}
assistance. The original derivations and reviews recorded effort
\texttt{ultra}; the 21 October manuscript integration of Notes 17--18
and associated discussion recorded effort \texttt{max}.
Derivations and reviews are
internal LLM work. Detailed proofs and corrections are in
\href{../notes/12_FINITE_EVALUATION_AND_ENDPOINT_INVERSE_COSTS_20261020.md}{Note 12},
\Href, \Nref, \CNref, \CTref, and
\href{../reviews/16_GEVREY_GLOBAL_APPROXIMATION_REVIEW_20261020.md}{Review 16},
\href{../reviews/17_NORMALIZED_HEAT_LIMIT_AND_ARITHMETIC_REVIEW_20261020.md}{Review 17}
and \href{../reviews/18_ARITHMETIC_HEAT_TRANSFER_AND_COMPLETION_REVIEW_20261021.md}{Review 18}.
\texttt{../DRAFT\_HISTORY.md} indexes milestones of this canonical draft.

\clearpage
\begin{thebibliography}{9}\small
\bibitem{Burnol}
J.-F. Burnol, \emph{A lower bound in an approximation problem involving
the zeros of the Riemann zeta function}, Theorem 1.3.
\href{https://arxiv.org/abs/math/0103058}{arXiv:math/0103058}.
\bibitem{CCM}
E. Carneiro, V. Chandee and M. B. Milinovich,
\emph{Bounding \(S(t)\) and \(S_1(t)\) on the Riemann hypothesis},
Theorem 2. \href{https://arxiv.org/abs/1309.1526}{arXiv:1309.1526}.
\bibitem{CS}
V. Chandee and K. Soundararajan,
\emph{Bounding \(|\zeta(1/2+it)|\) on the Riemann hypothesis},
Theorem 1.1. \href{https://arxiv.org/abs/0908.2008}{arXiv:0908.2008}.
\bibitem{Polymath}
D. H. J. Polymath, \emph{Effective approximation of heat flow evolution
of the Riemann $\xi$ function, and a new upper bound for the de
Bruijn--Newman constant}, Proposition 9.1(i).
\href{https://arxiv.org/abs/1904.12438v2}{arXiv:1904.12438v2}.
\bibitem{Redheffer}
R. M. Redheffer, \emph{On a theorem of Plancherel and P\'olya},
Pacific J. Math. \textbf{3} (1953), 823--835, Section 5, Theorem 11.
\href{https://msp.org/pjm/1953/3-4/pjm-v3-n4-p11-s.pdf}{Primary text}.
\end{thebibliography}
\end{document}
